% ===========================================================================
% Résumé / Abstract / Mots-clés.
%
% Chaque chiffre avancé ici est REPRIS et démontré plus loin dans le rapport.
% Rien ne doit apparaître dans un résumé qui ne soit établi au chapitre 5.
% ===========================================================================
\chapter*{Résumé}
\addcontentsline{toc}{chapter}{Résumé}

Ce rapport présente la conception et la réalisation de PIQBIT, une plateforme
de recrutement augmentée par l'apprentissage automatique. Le projet répond à
une difficulté concrète du traitement des candidatures : lorsqu'une offre
attire des dizaines de curriculum vitæ, leur lecture manuelle devient longue,
peu reproductible et difficile à justifier. Les outils de suivi de candidatures
classiques stockent les dossiers sans les comprendre ; les filtres par
mots-clés, eux, écartent des profils pertinents dès que le vocabulaire du
candidat s'éloigne de celui de l'annonce.

PIQBIT propose une chaîne de traitement complète et traçable. Le curriculum
vitæ déposé est analysé automatiquement — extraction du texte, des compétences,
de l'expérience et de la formation — puis noté par rapport aux exigences de
l'offre, et le candidat est classé. Un rapprochement sémantique par vecteurs
permet de retrouver des profils dont le vocabulaire diffère de celui de
l'annonce. La plateforme couvre ensuite tout le cycle de recrutement :
planification d'entretiens avec confirmation de créneau par le candidat et
synchronisation d'agenda, test technique chronométré, prédiction salariale par
un modèle entraîné, assistance à la négociation, édition et signature
électronique du contrat, puis activation de la fiche employé.

La solution est réalisée en Angular 19 côté client et FastAPI côté serveur,
sur PostgreSQL 16 doté d'une extension vectorielle, avec une politique
d'autorisation externalisée et un schéma piloté par migrations. Elle expose
173 routes réparties en 18 modules, s'appuie sur 33 services applicatifs et
25 tables.

Les principales contributions portent sur la notation explicable des
candidats, le rapprochement sémantique avec repli fonctionnel, l'assistant
extractif dont chaque réponse cite ses sources, et un modèle de prédiction
salariale sélectionné parmi huit candidats sur l'erreur de test — et non
d'entraînement — atteignant une erreur absolue moyenne de 402 dinars. Le
rapport expose également, sans les atténuer, les limites constatées : une
incohérence d'autorisation révélée par la campagne de validation, un budget de
taille de bundle non tenu, et l'absence d'environnement d'hébergement cible.

\chapter*{Abstract}
\addcontentsline{toc}{chapter}{Abstract}

This report presents the design and implementation of PIQBIT, a
machine-learning-assisted recruitment platform. The project addresses a
concrete difficulty in application processing: when a job opening attracts
dozens of curricula vitae, reading them manually becomes slow, hardly
reproducible, and difficult to justify. Conventional applicant-tracking systems
store documents without understanding them, while keyword filters discard
relevant candidates as soon as their vocabulary departs from the wording of the
advertisement.

PIQBIT provides a complete and traceable processing chain. An uploaded
curriculum vitæ is analysed automatically — text, skills, experience and
education extraction — then scored against the requirements of the job offer,
and the candidate is ranked. Vector-based semantic matching retrieves profiles
whose wording differs from the advertisement. The platform then covers the
whole recruitment cycle: interview scheduling with candidate-side slot
confirmation and calendar synchronisation, timed technical assessment, salary
prediction from a trained model, negotiation support, electronic contract
signature, and employee record activation.

The solution is built with Angular 19 on the client side and FastAPI on the
server side, over PostgreSQL 16 with a vector extension, an externalised
authorisation policy, and a migration-driven schema. It exposes 173 routes
across 18 modules, relies on 33 application services and 25 tables.

The main contributions are explainable candidate scoring, semantic matching
with a functional fallback, an extractive assistant whose answers cite their
sources, and a salary-prediction model selected among eight candidates on test
error — not training error — reaching a mean absolute error of 402 dinars. The
report also states, without softening them, the limitations observed: an
authorisation inconsistency revealed by the validation campaign, an unmet
bundle-size budget, and the absence of a target hosting environment.

\chapter*{Mots-clés}
\addcontentsline{toc}{chapter}{Mots-clés}

\begin{multicols}{2}
\begin{itemize}[leftmargin=*,itemsep=4pt]
  \item Recrutement assisté
  \item Analyse automatique de CV
  \item Notation de candidats
  \item Recherche sémantique
  \item Plongements lexicaux
  \item Prédiction salariale
  \item Signature électronique
  \item Contrôle d'accès par rôle
  \item FastAPI
  \item Angular
\end{itemize}
\end{multicols}
